\chapter{Analisis dan Perancangan}
\emph{Cara menulis bab ini.} Bab Analisis dan Perancangan menyajikan analisis kebutuhan dan rancangan sistem (proses, basis data, dan antarmuka).


\section{Analisis Kebutuhan}
\emph{Cara menulis.} Identifikasi kebutuhan fungsional dan non-fungsional berdasarkan rumusan masalah.

Analisis kebutuhan dibagi menjadi kebutuhan fungsional dan nonfungsional.
Kebutuhan fungsional sistem dirangkum pada Tabel~\ref{tab:fungsional}.

\begin{table}[H]
\centering
\caption{Kebutuhan fungsional sistem}
\label{tab:fungsional}
\begin{tabular}{|c|l|}
\hline
\textbf{Kode} & \textbf{Kebutuhan Fungsional} \\ \hline
\multicolumn{1}{|c|}{(1)} & \multicolumn{1}{c|}{(2)} \\ \hline
KF-01 & Sistem dapat melakukan autentikasi pengguna. \\
KF-02 & Sistem dapat mengunggah dan memvalidasi data. \\
KF-03 & Sistem dapat menampilkan hasil prediksi. \\ \hline
\end{tabular}
\end{table}

\section{Perancangan Proses (Use Case)}
\emph{Cara menulis.} Gambarkan proses dengan \emph{use case}/diagram alir beserta deskripsinya.

Interaksi pengguna dengan sistem digambarkan pada \emph{use case diagram}
Gambar~\ref{gbr:usecase}.

\begin{figure}[H]
\centering
\begin{tikzpicture}
  \node (sys) [uc-sistem, minimum width=5cm, minimum height=4cm] at (2.2,0) {};
  \node[above] at (sys.north) {\small Sistem};
  \pic (akt) at (-2.2,0) {uc-aktor};
  \node[below] at (-2.2,-1.2) {\small Pengguna};
  \node[uc-case] (u1) at (2.2,1.2) {Login};
  \node[uc-case] (u2) at (2.2,0)   {Unggah Data};
  \node[uc-case] (u3) at (2.2,-1.2){Lihat Prediksi};
  \draw[uc-asosiasi] (-2.0,0) -- (u1);
  \draw[uc-asosiasi] (-2.0,0) -- (u2);
  \draw[uc-asosiasi] (-2.0,0) -- (u3);
\end{tikzpicture}
\caption{Use case diagram sistem}
\label{gbr:usecase}
\end{figure}

\section{Perancangan Basis Data}
\emph{Cara menulis.} Sajikan rancangan basis data (mis.\ ERD atau skema tabel).

Bagian ini memuat rancangan basis data (mis.\ \emph{Entity Relationship Diagram})
beserta struktur tabel.

\section{Perancangan Antarmuka}
\emph{Cara menulis.} Sajikan rancangan antarmuka (mis.\ \emph{wireframe} atau \emph{mockup}).

Bagian ini memuat rancangan antarmuka (\emph{wireframe}/mockup) halaman utama
sistem.
